\answer{ex:12:1}{$0.60$（若无相关则为$0.18$），极限$\sqrt{c/(rT)}\approx0.58$：一百个神经元只能区分“有/无”。}
\answer{ex:12:2}{(a)~$10^8$个脉冲，$\approx10\,\text{mJ}$：占全脑$3$\,J的$0.3\,\%$。(b)~$3$\,J，约为其$300$倍。}
\answer{ex:12:3}{(a)~$\theta_A\in[2,3)$，$\theta_B\in[1,2)$；另一个图形给出的输入分别不超过$2$和$1$。(b)~$10\cdot96^2\approx9.2\cdot10^4$。}
\answer{ex:12:4}{(a)~$T_d=0.85\,\tau_a=0.34\,\text{s}$。(b)~$T_d\propto\tau_a$，与$I$无关：$\tau_a\approx3.5\,\text{s}$（模拟：$3.1\,\text{s}$）。}
\answer{ex:12:5}{$p\approx0.17$。正确率下降缓慢：$N=8$、$12$、$24$、$48$、$96$、$192$时分别为$0.90$、$0.88$、$0.86$、$0.84$、$0.82$、$0.79$。}
\answer{ex:12:6}{$\tau_{\text{tr}}=20$、$50$和$200\ms$时十次运行全部无误（$30\ms$时十次中有一次失败）；$5$--$10\ms$时约为$0.5$，$0.5$--$2\,\text{s}$时质量降到$0.86$。上限：资格迹须在停顿期间消退，$e^{-0.3\,\text{s}/\tau_{\text{tr}}}\ll1$。}
\answer{ex:12:7}{一个延迟不够：$\pm5\ms$的窗口覆盖不了$5$到$20\ms$的$\Delta x/v$。需两条抑制支路，$5<d_1\le10\ms$，$15\le d_2\le d_1+10\ms$。}
\answer{ex:12:8}{\texttt{two\_stage}：$\Delta=2$，翻转一个点得到平局，翻转两个点才改变答案。数1的分类器：翻转一个点即可；$p=0.1$时为$0.57$。}
